\documentclass[11pt,a4paper]{article}
\usepackage[margin=2.2cm]{geometry}
% Preámbulo común de los documentos (article). Compatible con Overleaf (pdfLaTeX, sin -shell-escape).
\usepackage[utf8]{inputenc}
\usepackage[T1]{fontenc}
\usepackage[spanish,es-noshorthands]{babel}
\usepackage{lmodern}
\usepackage{microtype}
\usepackage{amsmath,amssymb}
\usepackage{booktabs}
\usepackage{array}
\usepackage{tabularx}
\usepackage{xcolor}
\usepackage{listings}
\usepackage{enumitem}
\usepackage{caption}
\usepackage[hidelinks]{hyperref}

\captionsetup{font=small,labelfont=bf}
\setlist{itemsep=1pt,topsep=3pt}

\lstset{
  basicstyle=\ttfamily\footnotesize,
  columns=fullflexible,
  keepspaces=true,
  frame=single,
  framerule=0.3pt,
  rulecolor=\color{gray!60},
  backgroundcolor=\color{gray!6},
  breaklines=true,
  inputencoding=utf8,
}

\newcommand{\rel}[1]{\textbf{#1}}
\newcommand{\code}[1]{\texttt{#1}}


\title{Diccionario de features\\[2pt]\large Listwise Learning to Rank sobre consultas LOINC}
\author{Machine Learned Ranking --- MUIA, UPM}
\date{Noviembre 2026}

\begin{document}
\maketitle

Cada par consulta--documento $(q,d)$ se representa con cinco features numéricas, calculadas en
\code{src/build\_dataset.py} \textbf{sin acceder a la etiqueta de relevancia}. El orden de las features es el
de los índices LETOR \code{1:}\ldots\code{5:}.

\section*{Preprocesado común}

\begin{itemize}
\item Texto en minúsculas y tokenización por secuencias alfanuméricas (\code{[a-z0-9]+}).
\item Eliminación de una lista corta de palabras vacías (\emph{in, of, or, by, the, for, and, from, to, a}).
\item Stemming mínimo del plural inglés regular (\emph{cells} $\rightarrow$ \emph{cell}, \emph{leukocytes} $\rightarrow$ \emph{leukocyte}).
\item El corpus de TF-IDF y BM25 es la colección de documentos \emph{únicos} de cada configuración: los 67 \code{long\_common\_name} del Excel para \code{train.txt} y los 1222 de la colección completa para \code{train\_terms.txt} y \code{train\_extended.txt}. No se usan las filas consulta--documento, que repetirían cada documento una vez por consulta.
\end{itemize}

\section*{Features}

\begin{center}
\small
\begin{tabularx}{\linewidth}{c >{\raggedright}p{3.1cm} >{\raggedright}X c}
\toprule
Índice & Nombre & Definición & Rango \\
\midrule
1 & \code{f1\_tfidf\_cosine} & Coseno entre los vectores TF-IDF (normalizados L2) de $q$ y de \code{long\_common\_name}. Es la feature del trabajo inicial, sin el límite \code{max\_features=20} y con la errata de la consulta corregida. & $[0,1]$ \\
2 & \code{f2\_bm25} & BM25 ($k_1{=}1.2$, $b{=}0.75$) de $q$ sobre \code{long\_common\_name}, dividido por el máximo de la consulta. & $[0,1]$ \\
3 & \code{f3\_synonym\_overlap} & Fracción de términos de $q$ presentes en el documento (\code{long\_common\_name}, \code{component}, \code{system}, \code{property}), aceptando sinónimos: \emph{white/cells} $\rightarrow$ \emph{leukocytes}, \emph{count} $\rightarrow$ \emph{\#/volume}, \emph{blood} $\rightarrow$ \emph{bld}, \emph{plasma} $\rightarrow$ \emph{plas}, \emph{serum} $\rightarrow$ \emph{ser}, \emph{platelets} $\rightarrow$ \emph{platelet/plt}. & $[0,1]$ \\
4 & \code{f4\_component\_match} & 1 si algún término de analito de la consulta aparece como palabra en \code{component}; 0 en otro caso. & $\{0,1\}$ \\
5 & \code{f5\_system\_match} & 1 si el término de muestra de la consulta (\emph{bld}, \emph{plas}, \emph{ser}) aparece literalmente en \code{system}; 0 en otro caso. Es binaria y no distingue compatible de incompatible. & $\{0,1\}$ \\
\bottomrule
\end{tabularx}
\end{center}

\section*{Fórmulas}

\[
f_1(q,d)=\frac{\mathbf{v}_q\cdot\mathbf{v}_d}{\lVert\mathbf{v}_q\rVert\,\lVert\mathbf{v}_d\rVert},\qquad
\mathrm{BM25}(q,d)=\sum_{t\in q}\mathrm{idf}(t)\,\frac{tf_{t,d}\,(k_1+1)}{tf_{t,d}+k_1\left(1-b+b\frac{|d|}{\mathrm{avgdl}}\right)},\qquad
f_2=\frac{\mathrm{BM25}(q,d)}{\max_{d'}\mathrm{BM25}(q,d')}
\]
con $\mathrm{idf}(t)=\ln\!\left(1+\frac{N-n_t+0.5}{n_t+0.5}\right)$, $N$ = número de documentos de la colección (67 o 1222).
\[
f_3(q,d)=\frac{1}{|T_q|}\sum_{t\in T_q}\mathbb{1}\big[\{t\}\cup\mathrm{syn}(t)\ \cap\ \mathrm{tokens}(d)\neq\emptyset\big]
\]

\section*{Por qué estas features y no otras}

\begin{itemize}
\item \textbf{Dependen de la consulta.} El trabajo inicial usaba \code{component}, \code{system} y \code{property}
  codificadas con \code{LabelEncoder}. Esos índices alfabéticos son iguales para cualquier consulta, y AdaRank usa cada feature
  como ranker débil, por lo que ordenar por ``posición alfabética del sistema'' no aporta información. Se han eliminado.
\item \textbf{Sin fuga de etiquetas.} Ninguna feature usa \code{relevance} ni el ranking ideal.
\item \textbf{Más gruesas que la regla de etiquetado.} $f_4$ y $f_5$ se inspiran en las mismas columnas que la regla, pero son
  binarias y no distinguen exacto/relacionado ni compatible/incompatible. Aun así, su parecido con la regla es una
  limitación que se declara en el informe.
\end{itemize}

\section*{Ejemplo (formato LETOR, \code{data/train.txt})}

% Generado por src/export_tables.py — no editar a mano
\begin{lstlisting}
3 qid:1 1:0.750959 2:1 3:1.01765 4:1 5:1 6:1 # docid = 74774-1
2 qid:1 1:0.499538 2:0.694848 3:0.465047 4:1 5:0 6:1 # docid = 14749-6
1 qid:1 1:0.428328 2:0.75537 3:0.5 4:1 5:0 6:1 # docid = 15076-3
3 qid:2 1:0.653337 2:1 3:1.03668 4:1 5:1 6:1 # docid = 1975-2
3 qid:3 1:0.467694 2:0.965348 3:1.01449 4:0.333333 5:1 6:1 # docid = 26464-8
1 qid:3 1:0.308181 2:0.881265 3:0.761091 4:0.333333 5:1 6:0 # docid = 26478-8
0 qid:1 1:0.237446 2:0.458793 3:0.491787 4:0 5:0 6:0 # docid = 933-2
\end{lstlisting}


Cada línea contiene la etiqueta, el \code{qid} (1 = glucose in blood, 2 = bilirubin in plasma, 3 = white blood cells count),
las features \code{índice:valor} y, en comentario, el código LOINC (\code{docid}). El archivo tiene 201 líneas, agrupadas en 67 por \code{qid}.

\begin{center}\small
\begin{tabular}{lll}
\toprule
Archivo & Contenido & Líneas \\
\midrule
\code{data/train.txt} & 3 consultas oficiales $\times$ 67 documentos del Excel (obligatorio) & 201 \\
\code{data/train\_terms.txt} & 3 consultas oficiales $\times$ 1222 documentos (más términos) & 3666 \\
\code{data/train\_extended.txt} & 15 consultas $\times$ 1222 documentos (más términos y consultas) & 18\,330 \\
\bottomrule
\end{tabular}
\end{center}

Los \code{qid} 4--15 son las consultas de la ampliación (selección estratificada, ver la guía de etiquetado): creatinine in serum,
red blood cells count, oxygen partial pressure in blood, prothrombin time in plasma, ketones in urine, SARS-CoV-2 RNA in respiratory
specimen, bacteria in urine culture, leukocyte esterase in urine, perinuclear ANCA in serum, HIV antibody and antigen in serum,
tacrolimus in blood y EGFR antigen in tissue. Sus sinónimos para $f_3$ (p.\,ej. \emph{red}$\to$\emph{erythrocyte},
\emph{prothrombin}$\to$\emph{PT/INR}, \emph{antibody}$\to$\emph{Ab}) están en \code{src/queries.py}.

\end{document}
